\documentclass[11pt,a4paper]{article}

\usepackage[margin=2.5cm]{geometry}
\usepackage{booktabs}
\usepackage{hyperref}
\usepackage{url}

\title{Milestone~1 Part~1 (M11): Parent-Position and Representation-Order Control\\[0.5em]
\large Supervisor-facing findings log}
\author{Causal ABA Learning project --- Samuel Waugh}
\date{18 June 2026 \quad (initial grid); 22 June 2026 \quad (status: \textbf{analysed} --- Stages 0--6 complete)}

\begin{document}
\maketitle

\begin{abstract}
Experiment~M11 (m1.1) tests whether target-wise ABA Learning on complete noiseless \emph{target-mechanism} tables---one isolated predictor plus a deterministic parent$\rightarrow$\texttt{x2} copy---recovers rules that track the \emph{direct-parent role} under variable-identity ($\pi$) and predictor-order ($\sigma$) transformations. This is \textbf{not} chain SCM sampling or causal discovery. \textbf{Answer:} parent-\emph{role} tracking holds for \textbf{binary} under fixed nd folding with full $\sigma$- and $\pi$-invariance; for \textbf{cat3} it holds \textbf{when the parent feature wins the first fold}, but $\sigma$-invariance \textbf{fails} under the default pipeline because BK block serialisation order fixes which value predicate \texttt{select\_rule} tries first, the entailment gate accepts or rejects that fold, and failed target entailment commits to assumptions without backtracking to the parent separator. Stage~6 ablations (ABL-100--107; commit \texttt{93d51bf}) confirm this chain: reordering BK on cat3~A recovers the singleton (ABL-101); \texttt{folding\_mode(all)} without BK change does not (ABL-105); binary remains $\sigma$-invariant (ABL-107). Part~1 is \textbf{analysed}; M1.2 (greedy vs nd) is unblocked as a comparator only.
\end{abstract}

\section{Position in Milestone~1}

Milestone~1 (see \texttt{docs/research/milestone\_plans/milestone1-plan.md}) investigates whether ABA Learning supports causal-discovery-adjacent parent-set recovery on controlled fixtures.

\begin{itemize}
  \item \textbf{Part~1 (m1.1, this document):} metamorphic validity control on the QL2 parent-recovery proxy. \textbf{Status: analysed} --- Stages 0--6 complete (grid, trace mechanism, ablations).
  \item \textbf{Part~2 (m1.2):} greedy vs nd on the m1.1 grid --- \textbf{next comparator experiment}; not required to explain cat3 $\sigma$ failure (mechanism established in Part~1).
\end{itemize}

Motivation: QI-002's categorical chain returned a rule citing \texttt{x0} where the expected direct parent was \texttt{x1}. m1.1 separates direct-parent-role tracking from fixed variable identity and from fixed representation position.

\section{Research question and design}

\subsection{Question}

Under complete, noiseless QL2-style target-mechanism tables, when the direct-parent and non-parent roles are exchanged between \texttt{x0} and \texttt{x1}, and the learner-visible predictor order is independently reversed, does the learnt \texttt{x2} rule follow the variable occupying the direct-parent role?

\subsection{Factorial structure}

Eight deterministic cells: two encodings (\{binary, cat3\}) $\times$ four cells \{A, B, C, D\}, derived from one canonical cell~A per encoding via:
\begin{itemize}
  \item $\sigma$: reverse learner-visible predictor order (values, examples, roles unchanged);
  \item $\pi$: exchange \texttt{x0}$\leftrightarrow$\texttt{x1} (names, values, edges relabelled; \texttt{x2} fixed).
\end{itemize}

Tabular DGP (all cells): exhaustive factorial over \((x0,x1)\); \texttt{x2 := parent} deterministically; the non-parent is marginally independent of \texttt{x2}. Each cell is one \emph{isolated} predictor plus a \emph{connected} parent--target pair, not a faithful chain sample.

\begin{center}
\small
\begin{tabular}{@{}llll@{}}
\toprule
Cells & Parent & Isolated & Data-level graph \\
\midrule
A, B & \texttt{x1} & \texttt{x0} & \texttt{x0} ; \texttt{x1}$\rightarrow$\texttt{x2} \\
C, D & \texttt{x0} & \texttt{x1} & \texttt{x1} ; \texttt{x0}$\rightarrow$\texttt{x2} \\
\bottomrule
\end{tabular}
\end{center}

Minimal complete factorial: binary 4 rows; cat3 9 rows; one row per \((x0,x1)\) assignment; no noise; seed~0.

\subsection{Pre-specified expected rules}

\begin{center}
\begin{tabular}{@{}lll@{}}
\toprule
Cell & Binary & Categorical-3 \\
\midrule
A, B & \texttt{x2(A) <- x1(A)} & \texttt{x2(A) <- x1\_val\_2(A)} \\
C, D & \texttt{x2(A) <- x0(A)} & \texttt{x2(A) <- x0\_val\_2(A)} \\
\bottomrule
\end{tabular}
\end{center}

These are the unique intended one-literal observable-feature separators in each cell, not uniqueness over the full ABA hypothesis space.

\subsection{Fixed learner settings}

Target \texttt{x2} only; default \texttt{folding\_mode: nd}; \texttt{folding\_steps: 15}; \texttt{prolog\_timeout\_s: 120}; cautious + relto (inherited QI-002-style). \textbf{Not} Russo-style Causal ABA.

\section{What was done}

\begin{enumerate}
  \item \textbf{Stage~0} --- Prolog-free fixture validation (\texttt{pytest}, \texttt{m11\_validate\_fixtures.py}).
  \item \textbf{Stage~1} --- eight-cell learning grid (\texttt{run\_grid}, \texttt{M11\_parent\_position.yaml}).
  \item \textbf{Stage~2} --- transformation-square verdicts ($\sigma$, $\pi$) at rule level.
  \item \textbf{Stage~3} --- triggered inspection of cat3 $\sigma$-pair A vs~B via \texttt{prolog.stdout}.
  \item \textbf{Stage~4} --- granular trace audit (cat3 A vs B): \texttt{read\_bk} $\rightarrow$ \texttt{fwt} $\rightarrow$ \texttt{select\_rule} $\rightarrow$ entailment vs assumption commitment.
  \item \textbf{Stage~5} --- runner/bridge clarity (\textbf{partial}): data-flow summary in Stage~4; full call-chain doc optional.
  \item \textbf{Stage~6} --- falsification-first ablations ABL-100--107 (\texttt{m11\_ablation\_run.py}, \texttt{m11\_ablation\_audit.py}).
\end{enumerate}

\subsection{Commands run}

\textbf{Grid (Phase~A):}
\begin{verbatim}
conda activate aba-asp
cd "/Users/samuelwaugh/Desktop/Causal ABA Learning/aba_asp"
python -m pytest causal/tests/test_m11_fixtures.py -q
python causal/scripts/m11_validate_fixtures.py
python -m causal.experiments.run_grid \
  --config causal/configs/experiments/M11_parent_position.yaml --no-resume
\end{verbatim}
Outcome: \texttt{completed=8}, all \texttt{solved}.

\textbf{Ablations (Phase~B, Stage~6):}
\begin{verbatim}
python -m causal.scripts.m11_ablation_run --all-required
python -m pytest causal/tests/test_m11_ablation_bk.py \
  causal/tests/test_m11_ablation_inspect.py -q
\end{verbatim}
Commit: \texttt{93d51bf}. Gating: ABL-100 audit\_pass $\rightarrow$ ABL-101 pass $\rightarrow$ ABL-102, 103, 105, 107; ABL-106 skipped (ABL-105 conclusive).

\subsection{Implementation}

\begin{itemize}
  \item \textbf{Fixtures:} \texttt{causal/experiments/handcrafted\_m11.py}.
  \item \textbf{Grid:} \texttt{causal/experiments/run\_grid.py}; config \texttt{M11\_parent\_position.yaml}.
  \item \textbf{BK order:} \texttt{causal/argcausaldisco\_integration.py} serialises by \texttt{DataFrame} column order (no sort).
  \item \textbf{Ablations:} \texttt{m11\_ablation\_bk.py}, \texttt{m11\_ablation\_inspect.py}, \texttt{M11\_ablations.yaml}, \texttt{m11\_ablation\_run.py}, \texttt{m11\_ablation\_audit.py}.
\end{itemize}

\section{Grid results (Phase~A)}

\subsection{Cell-level (Stage~1)}

\begin{center}
\small
\begin{tabular}{@{}llllp{3.2cm}p{3.5cm}l@{}}
\toprule
Enc. & Cell & Parent & Order & Expected & Actual (summary) & Class. \\
\midrule
binary & A & x1 & [x0,x1] & \texttt{x2<-x1} & \texttt{x2 :- x1.} & exact \\
binary & B & x1 & [x1,x0] & \texttt{x2<-x1} & \texttt{x2 :- x1.} & exact \\
binary & C & x0 & [x0,x1] & \texttt{x2<-x0} & \texttt{x2 :- x0.} & exact \\
binary & D & x0 & [x1,x0] & \texttt{x2<-x0} & \texttt{x2 :- x0.} & exact \\
cat3 & A & x1 & [x0,x1] & \texttt{x2<-x1\_val\_2} & 3-rule assumption; \texttt{x0\_val\_*} & discrepant$^\dagger$ \\
cat3 & B & x1 & [x1,x0] & \texttt{x2<-x1\_val\_2} & \texttt{x2 :- x1\_val\_2.} & exact \\
cat3 & C & x0 & [x0,x1] & \texttt{x2<-x0\_val\_2} & \texttt{x2 :- x0\_val\_2.} & exact \\
cat3 & D & x0 & [x1,x0] & \texttt{x2<-x0\_val\_2} & 3-rule assumption; \texttt{x1\_val\_*} & discrepant$^\dagger$ \\
\bottomrule
\end{tabular}
\end{center}

{\footnotesize $^\dagger$All eight cells are \texttt{solved}; discrepant = assumption-based rule structure vs pre-specified singleton.}

\subsection{Transformation-level (Stage~2)}

\begin{center}
\begin{tabular}{@{}lllll@{}}
\toprule
Enc. & Property & Pair & Rule-level & Role-level \\
\midrule
binary & $\sigma$ & A$\leftrightarrow$B & holds & holds \\
binary & $\sigma$ & D$\leftrightarrow$C & holds & holds \\
binary & $\pi$ & A$\leftrightarrow$D & holds & holds \\
binary & $\pi$ & B$\leftrightarrow$C & holds & holds \\
cat3 & $\sigma$ & A$\leftrightarrow$B & \textbf{fails} & \textbf{fails} \\
cat3 & $\sigma$ & D$\leftrightarrow$C & \textbf{fails} & \textbf{fails} \\
cat3 & $\pi$ & A$\leftrightarrow$D & \textbf{fails} & \textbf{fails} \\
cat3 & $\pi$ & B$\leftrightarrow$C & holds & holds \\
\bottomrule
\end{tabular}
\end{center}

Cat3 $\pi$ on A$\leftrightarrow$D fails at expected-rule level: both return assumption structures (symmetric failure class), not the pre-specified singleton.

\subsection{Secondary metrics}

\begin{center}
\begin{tabular}{@{}lrrrrr@{}}
\toprule
Cell & clean\_rec. & $n_{\mathrm{rules}}$ & $n_{\mathrm{asm}}$ & wall (s) & cov\_py; cov\_pl \\
\midrule
binary A--D & 1 & 1 & 0 & 0.3--0.4 & (1,1);(1,1) each \\
cat3 B, C & 1 & 1 & 0 & 0.3--0.4 & (1,1);(1,1) each \\
cat3 A & 0 & 3 & 3 & 1.54 & (1,0);(0,1) \\
cat3 D & 0 & 3 & 3 & 1.55 & (1,0);(0,1) \\
\bottomrule
\end{tabular}
\end{center}

Classifications use \texttt{bk.sol.aba} and \texttt{prolog.stdout}, not \texttt{cov\_py} alone.

\section{Mechanistic account (Stage~4)}

\textbf{Observation (Stage~3):} cat3 A vs B diverge at the first nd fold on rote positive \texttt{x2(A)<-[A=3]} (trace line~35).

\textbf{Inputs:} same E+, E-, and BK rule multiset; $\sigma$ differs only in BK block serialisation order.

\textbf{Engine path:} \texttt{read\_bk} assigns rule IDs 1--18 in file order; \texttt{fwt} indexes equalities \texttt{A=n} under \texttt{=/2}; \texttt{select\_rule} picks the first subsuming match for \texttt{[A=3]} (ID~3 = third BK line).

\begin{center}
\begin{tabular}{@{}llll@{}}
\toprule
Cell & ID~3 rule & Folded target rule & Entailment \\
\midrule
cat3 A & \texttt{x0\_val\_0(A):-A=3} & \texttt{x2(A) <- [x0\_val\_0(A)]} & \textbf{fails} (covers E$-$ rows 1,2) \\
cat3 B & \texttt{x1\_val\_2(A):-A=3} & \texttt{x2(A) <- [x1\_val\_2(A)]} & \textbf{passes} (E$+$ only) \\
\bottomrule
\end{tabular}
\end{center}

\textbf{Why A cannot recover B's singleton within the same run:}
\begin{enumerate}
  \item Failed \texttt{gen2} on the first target fold does \textbf{not} backtrack \texttt{select\_rule} to the parent rule at ID~12; engine commits to new assumption \texttt{alpha\_1}.
  \item Rote \texttt{x2(A)<-[A=3]} is consumed; assumption-augmented \texttt{x2} rules are not re-folded (\texttt{nonintensional/1}).
  \item Later \texttt{x1\_val\_*} folds apply to \texttt{c\_alpha\_*} contrary heads only --- exception repair, not parent singleton recovery.
\end{enumerate}

\textbf{Trace shape:} A: 173 lines, 27 ABA rules; B: 54 lines, 19 rules. Cat3 D vs C mirrors A vs B under $\pi$ (confirmed ABL-100, ABL-103).

\textbf{Ablation confirmation:} ABL-101 on cat3~A config (BK reordered to B block order, same data/examples) reproduces B's trace shape: 53 lines, 19 rules, 0 assumptions, singleton \texttt{x2(A):-x1\_val\_2(A)}.

Code: \texttt{asp\_utils.pl} (\texttt{read\_bk}, \texttt{update\_fwt}); \texttt{folding.pl} (\texttt{select\_rule}); \texttt{gen.pl} (\texttt{gen1}, \texttt{gen2}, \texttt{gen6}).

\section{Ablation evidence (Stage~6)}

Pipeline artefacts: \texttt{causal/outputs/aba\_learning/grid/M11\_ablations/<ABL-id>/} (\texttt{abl\_summary.json}, \texttt{prolog.stdout}, \texttt{bk.sol.aba}).

\subsection{Observables summary}

\begin{center}
\small
\begin{tabular}{@{}llllrrl@{}}
\toprule
ID & Intervention & first\_fold & entail. & trace & rules & Target outcome & Verdict \\
\midrule
100 & audit C vs D & \texttt{x0\_val\_2}/\texttt{x1\_val\_0} & pass/fail & 53/172 & 19/27 & singleton / assumptions & audit\_pass \\
101 & BK order A$\leftarrow$B & \texttt{x1\_val\_2} & pass & 53 & 19 & singleton & pass \\
102 & prepend parent rule & \texttt{x1\_val\_2} & pass & 53 & 20 & singleton (+1 BK rule) & pass \\
103 & BK order C$\leftarrow$D & \texttt{x1\_val\_0} & fail & 172 & 27 & assumption triple & pass \\
104 & audit A vs B & \texttt{x0\_val\_0}/\texttt{x1\_val\_2} & fail/pass & 172/53 & 27/19 & grid refs & audit\_pass \\
105 & \texttt{folding\_mode(all)} on A & \texttt{x0\_val\_0} & fail & 157 & 27 & assumption triple & pass \\
106 & \texttt{asm\_intro(sechk)} & --- & --- & --- & --- & skipped & --- \\
107 & audit binary A vs B & \texttt{x1}/\texttt{x1} & pass/pass & 50/50 & 5/5 & $\sigma$-invariant & audit\_pass \\
\bottomrule
\end{tabular}
\end{center}

ABL-106 skipped: ABL-105 conclusive for H3.

\subsection{Per-hypothesis verdicts}

\begin{description}
  \item[H$\sigma\leftrightarrow$BK / H1 (order $\rightarrow$ first fold).] ABL-101: reordering BK blocks on cat3~A (rule multiset preserved) changes first fold from \texttt{x0\_val\_0} to \texttt{x1\_val\_2} and outcome from assumptions to singleton. ABL-102: inserting \texttt{x1\_val\_2(A):-A=3} as rule ID~1 alone suffices --- lowest matching ID wins, not block narrative alone. \textbf{Supported.}
  \item[H2 (fold semantics $\rightarrow$ entailment).] ABL-104 (grid A vs B): isolated-predictor fold fails entailment; parent separator passes. ABL-101/102: parent-first fold passes. \textbf{Supported.}
  \item[H3 (failed entailment $\rightarrow$ no recovery).] ABL-105: \texttt{folding\_mode(all)} on unchanged cat3~A BK still yields assumption triple (157 lines, 27 rules); contrast ABL-101 singleton without engine-mode change. \textbf{Supported.}
  \item[H$\pi$ (symmetry under parent swap).] ABL-100: C singleton / D assumptions mirrors B/A under $\pi$. ABL-103: reversing BK on C yields D-like failure. \textbf{Supported.}
  \item[Hbin (encoding moderator).] ABL-107: binary A and B both fold parent \texttt{x1} first with pass entailment --- $\sigma$-invariant. Cat3 fails $\sigma$ because value predicates create competing subsume matches disambiguated by BK ID order. \textbf{Supported.}
\end{description}

\subsection{What ablations establish vs Stage~4 alone}

Stage~4 correlated $\sigma$ (BK order) with outcome via trace inspection. Ablations show BK order is \textbf{sufficient} to flip outcome on fixed data/examples (101, 102) and that engine folding mode alone is \textbf{not sufficient} to recover singleton (105). Interventions are post-hoc BK edits; they do not modify the default bridge serialisation.

\section{Integrated conclusion}

\subsection{Answer to the research question}

\begin{center}
\begin{tabular}{@{}llll@{}}
\toprule
Encoding & Parent-role tracking & $\sigma$-invariance (default pipeline) & Explanation \\
\midrule
binary & Yes (all four cells) & Yes & Parent wins first fold under both BK orders (ABL-107) \\
cat3 & Yes \textbf{when parent feature wins first fold} & \textbf{No} & Competing value predicates + BK ID order; non-parent-first $\rightarrow$ assumptions \\
\bottomrule
\end{tabular}
\end{center}

The learner \textbf{can} track the direct-parent role in both encodings. Under default nd + cautious + relto, cat3 $\sigma$-invariance fails because representation order permutes BK serialisation, which fixes \texttt{select\_rule}'s first match, which gates entailment, which commits or avoids the assumption path --- not because examples or data omit the parent.

\subsection{Causal chain (ablation-supported, bounded)}

\[
\sigma \;\Rightarrow\; \text{BK block order} \;\Rightarrow\; \text{rule ID at first match}
\;\Rightarrow\; \text{first-fold predicate} \;\Rightarrow\; \text{entailment pass/fail}
\;\Rightarrow\; \text{singleton vs assumption commitment.}
\]

This is an \textbf{ABA Learning input-serialisation / search-control effect} under fixed nd + relto. It is \textbf{not} causal discovery, not Russo-style Causal ABA, and not a claim that the engine is globally ``name biased'' --- the same parent rule is recovered when BK order favours it (ABL-101).

\subsection{Explicit non-claims}

\begin{itemize}
  \item Not chain SCM sampling, chain discovery, or whole-graph recovery.
  \item Not Russo-style Causal ABA.
  \item Not a claim that greedy folding fixes cat3 $\sigma$ failure (M1.2 comparator only).
  \item Not a test of correlated-ancestor / faithful chain SCM behaviour (QI-003/QI-004).
  \item Ablation BK reorder is diagnostic, not a production fix to \texttt{argcausaldisco\_integration.py}.
\end{itemize}

\section{Supervisor guidance: origin and resolution (June 2026)}

Fabrizio's review (\texttt{docs/research/supervisor\_guidance.md}) reopened m1.1 after the initial eight-cell grid with these requirements:

\begin{itemize}
  \item Granular \texttt{prolog.stdout} inspection for cat3 pairs --- \textbf{addressed} (Stage~4; D vs C via ABL-100/103).
  \item State \textbf{why, when, and how} --- \textbf{addressed} (ID~3 table; entailment gate; assumption commitment; ablation chain).
  \item Ablation only after a formed hypothesis --- \textbf{addressed} (Stage~6 falsification-first plan executed).
  \item ND vs greedy not sufficient alone --- \textbf{addressed} (ABL-105; greedy remains M1.2 scope).
\end{itemize}

Part~1 scope expanded from verdict table to mechanistic account; that account is now \textbf{ablation-supported}. Status: \textbf{analysed}.

\section{Limitations}

\begin{itemize}
  \item Single seed; no noise; target-mechanism tables only; target-wise learning on \texttt{x2}.
  \item cat3 A/D are \texttt{solved} but discrepant vs pre-specified singleton criterion.
  \item Coverage metrics disagree on cat3 A/D assumption rules.
  \item Ablation interventions are post-hoc BK edits, not automatic pipeline fixes.
  \item Stage~5 full runner call-chain doc optional.
  \item Greedy smoke test on m1.1 grid informal (formal M1.2 grid not yet run).
\end{itemize}

\section{Next steps}

\begin{itemize}
  \item \textbf{M1.2:} rerun m1.1 eight-cell grid under \texttt{folding\_mode: greedy} as comparator (not blocker for Part~1 closure).
  \item Optional Stage~5 runner/bridge call-chain documentation.
  \item Milestone~2+: Causal-ABA-guided bridge (out of m1.1 scope).
\end{itemize}

\section{Artefact locations}

\textbf{Stage-0:} \texttt{causal/outputs/m11\_parent\_position/validation/}.

\textbf{Grid:} \texttt{causal/outputs/aba\_learning/grid/M11\_parent\_position/cells/m11\_<enc>\_<cell>/}.

\textbf{Ablations:} \texttt{causal/outputs/aba\_learning/grid/M11\_ablations/}; summary \texttt{summary/abl\_results.md}.

\textbf{Records:} \texttt{docs/experiments/qualitative/M1.1-parent-position.md}; ablation plan \texttt{milestone1\_part1\_ablations.md}.

\section{Claims cross-reference}

Supported claims in \texttt{docs/report/claims\_ledger.md}: M11-C-001 (binary control), M11-C-002 (cat3 $\sigma$ failure), M11-C-003 (order $\rightarrow$ first fold; ablation-supported), M11-C-005 (grid + closure), M11-C-006 (no recovery without BK reorder; ABL-105), M11-C-007 (integrated cat3 account). M11-C-004 (greedy fixes $\sigma$) weakened / pending M1.2.

\section{Project references}

\begin{itemize}
  \item Experiment record: \texttt{docs/experiments/qualitative/M1.1-parent-position.md}
  \item Milestone plan Part~1: \texttt{docs/research/milestone\_plans/milestone1\_part1\_parent\_position.md}
  \item Ablation plan: \texttt{docs/research/milestone\_plans/milestone1\_part1\_ablations.md}
  \item Milestone overview: \texttt{docs/research/milestone\_plans/milestone1-plan.md}
  \item Configs: \texttt{M11\_parent\_position.yaml}, \texttt{M11\_ablations.yaml}
  \item QI-002 baseline: \texttt{docs/experiments/qualitative/QI-002.md}
\end{itemize}

\end{document}
